\section{Signed residuals expose nearby disc sectors}
\label{sec:sector-windows}

The preceding planar work improves supplied-sector queries. The next step
connects those queries to actual diagram exteriors: after coherent candidates
fail, search compatible type assignments near one of their supports.
This produces a new source-certified disc for a known coherent miss.
The search remains optional because nearby-sector coverage is not guaranteed
and its cost on misses can exceed the incumbent fallback.

\subsection{A real disc outside the tested coherent families}

The three-crossing closure of the two-strand word $(1,1,-1)$ is an unknot.
Its canonical diagram exterior has sixty tetrahedra. The raw, tree and
minimum-span coherent candidates do not supply a disc or cappable annulus.
Complete standard enumeration of their actual sectors also misses a disc;
positive Euler rays in some of those sectors are inessential.

A radius-one type probe of the raw and optimized supports checks 360
distinct neighbouring assignments without success. At radius two about
the optimized support, an unfiltered scan finds a disc on its 4,920th
query, after 36.464 seconds. The two edits permit type one at tetrahedron
10 and type two at tetrahedron 12, both previously unoccupied. The resulting
sector has matching nullity two. Its first tested standard ray is a
primitive disc with every coordinate zero or one; only eight quadrilateral
types are actually used, although the allowed sector retains the other
seed types. Independent diagram-source and normal-disc replay accepts it.
The source-disc proof is independent of the unsuccessful coherent surface.

The two new columns cannot occur separately. Their joint exposure suggests
an exact filter based on matching equations rather than attempting every
nearby assignment. The filter reduces this example to two single-column
candidates and 73 pairs. The standalone filtered probe finds the same disc
after 31 queries in about 0.411 seconds, including residual planning.
These exploratory measurements motivate the implementation; the paired
native measurements below have separate, disclosed scopes.

\subsection{Eliminating triangles once for all type alternatives}

Use the actual standard matching equations, written
$B\tau+Dq=0$. Here $q$ has all $3t$ labelled quadrilateral coordinates;
the algebraic system temporarily permits all types so it can describe
alternatives. Admissibility is imposed again when a compatible sector is
queried. The triangle block is a signed graph-incidence transpose.
A spanning forest expresses corner triangles as integral linear potentials
in $q$, up to their free link offsets. Every non-tree edge contributes its
cycle equation. Pure-Q equations are also retained. The resulting sparse
matrix $A$ satisfies
\[
 \exists\tau\ (B\tau+Dq=0)\quad\Longleftrightarrow\quad Aq=0.
\]
This is exact linear elimination; no mod-two reduction or expanded normal
pieces are used. For nonnegative compatible $q$, subtracting minima and
adding link offsets gives the usual nonnegative canonical lift.

Let $S$ be the positive quadrilateral support of an actual validated seed
vector. Compute an exact column-space echelon basis for $A_S$, and let
$\pi$ be its linear quotient projection. For any nonbase type $j$, define
the residual $r_j=\pi(A_j)$. Thus $r_j=0$ precisely when $A_j$ is in the
old column span. Nonzero rational residuals are reduced to a primitive
integer direction with positive first entry, retaining a separate sign.
Equality of unsigned keys then means proportionality; opposite signs mean
a positive cancellation is possible. The scalar magnitudes need not agree.

The forest construction and sparse rational column elimination take
polynomial bit work and storage in the actual encoded source. Fundamental
cycle entries are bounded integers; elimination coefficients and reduced
residual keys have polynomial bit length by determinant bounds. Dictionary
collisions can add polynomial work but do not change this complexity class.
The residual matrix is a producer filter, not a certificate that the supplied
source is the exterior of a particular diagram.

\subsection{A complete radius-two filter, including replacements}

Represent a compatible type assignment by a word of length $t$ in
$\{\varnothing,0,1,2\}$. Distance counts tetrahedra whose assignments differ.
Deleting a type only sets its coordinate to zero and therefore restricts
to a coordinate face. It cannot create a standard extreme ray absent from
the containing cone. Other edits add a type at an empty site or replace
the previously allowed type.

\begin{proposition}[Signed residual window filter]
The union of non-link standard extreme rays over the full radius-two
assignment window about $S$ is preserved by querying only:
\begin{enumerate}
\item the base sector;
\item each single nonbase type with zero residual;
\item each pair at distinct tetrahedra with both residuals zero;
\item each pair at distinct tetrahedra with nonzero oppositely signed
      proportional residuals.
\end{enumerate}
For a radius-one window, only the first two kinds are needed. Replacements
remove their old type before adding the new one, preserving admissibility.
\end{proposition}
\begin{proof}
For any admitted normal vector, let $J$ be its positive nonbase types.
Within the window $|J|\le2$, and projection of its matching equation gives
\[
                  \sum_{j\in J}q_jr_j=0,\qquad q_j>0.
\]
If $J$ is empty its coordinates use only old types. If $|J|=1$, its
residual must vanish. If $|J|=2$, either both vanish or they are nonzero
opposite multiples. A pair with exactly one zero residual cannot cancel.
This proves that every possible positive nonbase support is on the list.

Unused edits need not be kept. If a queried assignment allows a new type
whose coordinate on a particular ray is zero, setting all quadrilaterals
at that site to zero gives a coordinate face containing the ray. This
face also belongs to the retained assignment that restores the old type
there. Extremality passes from the original cone to its face and then
from that face to the containing retained cone. Pure deletions are handled
the same way. Thus every original window ray remains extreme in one of the
retained sectors, rather than merely belonging to their feasible union.
Conversely every retained assignment is itself within the radius-two
window. The ray unions are therefore equal.
\end{proof}

There is no necessity-only positivity shortcut hidden in the actual
disc test: each retained sector still receives complete standard-ray
enumeration and the established exact disc observer. Opposite residuals
alone need not make a replacement feasible after its old columns are
removed. For additions at previously empty sites, the seed's strictly
positive old support does give sufficiency for feasibility: solve the
old-span equation, then add a sufficiently large multiple of that positive
seed to remove negative old coordinates. Feasibility still does not
establish disc topology or essentiality.

Writing $z$ for the number of zero residual columns and $p_h,m_h$ for
the signed bucket sizes in nonzero projective class $h$, the unpruned
candidate bound is
\[
                 1+z+\binom z2+\sum_h p_hm_h.
\]
Same-tetrahedron pairs are removed. This can be far below the naive
$1+3t+9\binom t2$ assignment count, but its worst case remains quadratic.
Nonzero residuals with the same sign are not paired; forgetting the sign
would waste precisely the positivity information used by this theorem.

\subsection{The logarithmic-radius conditional accounting}

The type-window perspective supplies a precise remaining structural
parameter. For a centre with matching nullity $d_0$, an edit deletes at
most one old column and adds at most one new column. Deletion cannot
increase kernel dimension, and addition increases it by at most one.
Every radius-$r$ sector therefore has nullity at most $d_0+r$.
The full assignment window has at most
\[
                   \sum_{j=0}^{r}\binom tj3^j
\]
members. By Section~\ref{sec:planar-adaptive}, complete local standard
enumeration consequently has total bit cost
\[
 \operatorname{poly}(t+B)(1+t)^{O(d_0+r)}.
\]
In particular, polynomial source size and $d_0,r=O(\log n)$ permit
quasi-polynomial complete window enumeration. The implemented residual
planner treats radii one and two; larger radii require positive dependencies
of several residuals rather than the pair classification above.

This is a supplied-window theorem. A general recognizer would still need
a bounded complete set of centres, a logarithmic-distance coverage theorem
for an essential vertex disc, polynomial essentiality testing and controlled
source construction. Finding one real disc at radius two proves none of
those universal claims. The earlier failed radius-one searches are retained
as counterevidence to simply activating the extra queries everywhere.

\subsection{Native integration and independent source binding}

\begin{sloppypar}
The optional \code{sector\_radius} argument of \code{normal\_seed\_decide}
accepts zero, one or two. Zero follows the prior code path exactly.
After all configured coherent attempts fail, the new stage chooses the
candidate with best Euler characteristic and then least piece span.
\code{sector\_residual.py} tests its base sector, prepares signed residual
groups, then queries single zero columns, opposite pairs and zero pairs.
Every stage shares the original work callback. Remaining orbit cycles
are shared across all window component queries and any earlier planar
capping queries. Budget interruption or incomplete component work cannot
become a negative knot claim.
\end{sloppypar}

A success uses the existing strict \code{diagram-normal-disc-v1} schema:
canonical input PD, actual exterior triangulation, normal coordinates and
the independently checked disc-count proof. The final checker authenticates
the diagram exterior and its disc component; it does not run residual
planning, coherent extraction or sector discovery. It neither requires nor
infers that this new disc is coherent. Boundary-shelling wrappers can now
carry this ordinary disc schema as well: the checker first binds the
original canonical exterior and every shelling move to the final source.
Callback errors remain observable even when inner geometry validators
convert malformed source errors to rejection.

\begin{sloppypar}
The recognition switch is \code{normal\_seed\_sector\_radius}; its CLI
counterpart is \code{--normal-seed-sector-radius}, used with the native
seed stage enabled. A positive result reports
\code{native-normal-sector-window}. The default remains zero, preserving
the earlier fast path. This optional source-certified accelerator coexists
with the complete recognition fallback; it supplies no negative verdict
from a missed local window.
\end{sloppypar}

\subsection{Measured coverage, costs and publication checks}
All 1,401 maintained tests pass in 234.317 seconds. The 26-test focused
run passes in 4.368 seconds, including actual CLI recognition. Literal
small-window ray unions match the residual-filtered union for radii one
and two. Independent dense matching confirms the base nullities. Tests
cover the new disc, missing source fields, altered coordinates, a different
diagram, work interruption, zero orbit allowance, false-valued cancellation
and the ordinary-disc schema transported through genuine boundary shellings.

The native audit uses the same 85-source corpus and verifies exact disabled
legacy equality on every input, including work, stages and proof contents.
Enabled positives increase from 24 to 31. The seven added source cases
have six distinct canonical PD codes, at crossing counts three through six;
they are not seven distinct knot types or a random sample of inputs.
Every added ordinary disc proof passes independent diagram-source replay
and a fresh Regina compressing-disc check. Six have coordinate maximum
one and the seventh maximum two. All seven use zero interval-orbit cycles
through the existing unit-ray certificate path.

Thirty-four audit cases hit the two-million shared work cap. Twenty finish
their requested window without finding a disc; seven stop at a new positive
and the 24 incumbent positives return before window construction. These
four counts account for all 85 cases. Capped cases do not prove absence,
and even complete local-window misses remain inconclusive for the knot.
The whole audit takes about 684.44 seconds, illustrating why default
activation is not justified by the additional witnesses alone. All 592
captured source pins remain fixed during audit and timing.

The exploratory evidence also retains the 2,340 distinct radius-one queries
about selected raw/optimized centres of four diagrams. Every query completes,
with matching nullity at most three and at most five emitted rays; none
finds a disc. These are genuine failures of that small neighbourhood, not
failures of the complete recognition fallback. The radius-two motivating
scan is a successful prefix, not an exhaustive negative experiment.
Its source, edits, coordinates, original scripts and proof are retained.

Five serial interleaved four-arm rounds complete 200 measured calls and
40 warmups. The seed-query times include all native construction, planning,
normal queries, successful source verification and serialized result output.
All configured candidate work completes in these uncapped small cases.
The separate complete-recognition rows include the established fallback
when the native stage misses; every arm reaches the expected final verdict.
All A/A arms, deterministic seed-result digests and full evidence byte sizes
are retained. These timings compare radius zero and radius two in the
same native implementation; the disabled implementation separately matches
the frozen \code{6c35c3e53} source on the complete audit.

\begin{samepage}
\noindent\textbf{Complete configured native seed queries and output.}
\begin{center}
\small
\begin{tabular}{lrrrrr}
\toprule
Input & off ms & on ms & off/on & off A/A & on A/A\\
\midrule
empty circle & 5.772 & 5.696 & 1.017 & 1.013 & 0.992\\
optimized positive & 57.391 & 67.196 & 0.858 & 1.004 & 1.186\\
genus one miss & 17.938 & 427.069 & 0.042 & 0.999 & 0.989\\
trefoil & 15.682 & 673.458 & 0.024 & 0.995 & 1.008\\
figure eight & 27.397 & 1402.391 & 0.019 & 1.001 & 0.994\\
\bottomrule
\end{tabular}
\end{center}

\end{samepage}

The genus-one example now produces a native disc, but its native query takes
about 427 milliseconds against the incumbent miss's 18 milliseconds.
That comparison is a coverage/cost tradeoff: the older native query returns
inconclusive and has not recognized the diagram. Early positive controls
avoid window construction. The optimized-positive seed row has a noisy
first-run ratio 0.858; a separate fifteen-round repeat gives ratio 1.005
with off/on A/A medians 1.010 and 0.989. Both runs are kept.

\begin{samepage}
\noindent\textbf{Complete forced-portfolio recognition.}
\begin{center}
\small
\begin{tabular}{lrrrrr}
\toprule
Input & off ms & on ms & off/on & off A/A & on A/A\\
\midrule
empty circle & 0.035 & 0.033 & 1.048 & 0.931 & 1.262\\
optimized positive & 50.211 & 51.131 & 1.007 & 1.008 & 1.005\\
genus one miss & 20.640 & 440.936 & 0.047 & 0.999 & 1.011\\
trefoil & 18.122 & 715.580 & 0.026 & 1.022 & 1.019\\
figure eight & 45.882 & 1503.935 & 0.031 & 1.454 & 0.997\\
\bottomrule
\end{tabular}
\end{center}

\end{samepage}

Complete recognition of the genus-one example takes about 441 milliseconds
with the new stage, against about 21 milliseconds with the incumbent's
complete fallback. The trefoil and figure-eight controls are likewise much
slower with extra windows. This is a regression, not a speedup in complete
recognition. A longer noisy-row repeat completes 180 measured calls and
twelve warmups. Its empty-circle control is at parity, and the figure-eight
off/on ratio is 0.022 with both A/A medians about 1.001. This confirms the
large regression rather than discarding it as first-run noise.

The stage therefore stays off by default. Its value is the additional
independently replayable normal witnesses and the polynomial residual filter
in place of an unfiltered local scan. No performance claim is made for
enabling it on every diagram. A future scheduler would need evidence that
its predicted window work is worthwhile against the available complete
fallback, in addition to any structural completeness theorem.

\begin{sloppypar}
Reproduction uses \code{normal\_orbit\_research.sector\_windows}, with
\code{audit --fresh-regina} or \code{benchmark --rounds 5}, from \code{fast/}.
Raw records, exploratory executions, test logs and the longer control run
are in \code{data/sector-windows-*}. The native runtime revision is
\code{a498521ef}. Publication checks independently replay all 31 enabled
positive proofs with coherent/sector/residual/orbit producers disabled,
verify frozen source pins and retain rendered-page review. The seven fresh
Regina observations support the actual new witnesses; the general window
and nullity statements depend on the proofs above.
\end{sloppypar}

